\section{Обзор существующих методов}
\subsection{Методы неявного представления поверхностей}

Современные подходы к представлению и хранению трехмерных сцен на основе знаковых полей расстояний (Signed Distance Fields, SDF) прошли значительную эволюцию: от классических неявных функций до сложных гибридных структур, ориентированных на ускорение вычислений на графических процессорах (Graphics Processing Unit, GPU).

\subsubsection{Классические аналитические представления} \label{sssec:analitical}

\subsubsection{Классические нейронные представления} \label{sssec:mlp}
Основополагающие методы, такие как NeuS \cite{wang2021neus} и VolSDF \cite{yariv2021volume}, используют полносвязные нейронные сети (Multilayer Perceptron, MLP) \cite{rumelhart1986learning} для параметризации SDF. В рамках данного подхода координата пространства напрямую отображается в значение расстояния и цвет с помощью объемного рендеринга \ref{sssec:vrender}. Несмотря на высокую геометрическую точность, данные методы характеризуются крайне низким темпом обучения (до 8--12 часов) и сложностями при рендеринга в реальном времени из-за необходимости многократного обращения к тяжелой нейронной сети для каждой пробы (sample) вдоль луча.

\subsubsection{Методы на основе хэш-сеток}
Для преодоления вычислительной избыточности MLP были предложены гибридные структуры данных. Стандартом де-факто стало использование многомасштабного хэш-кодирования \cite{muller2022instant} (Multiresolution Hash Encoding), популяризированного в работах Neuralangelo \cite{li2023neuralangelo} и NeuS2 \cite{wang2023neus2}. Данные методы используют хэш-таблицы как эффективные кодировщики признаков, что позволяет сократить размер MLP и ускорить процесс обучения в 100 и более раз по сравнению с чисто нейронными аналогами \cite{wang2021neus}\cite{yariv2021volume}.

\subsubsection{Разреженные воксельные структуры}
С точки зрения ускорения рендеринга на GPU, ключевой интерес представляют явные и иерархические структуры хранения:
\begin{itemize}
    \item \textbf{Разреженные Воксельные Решетки (Sparse Voxel Grids):} Методы типа GSLD \cite{yu2024gsdf3dgsmeetssdf} отказываются от нейронных сетей в пользу прямого хранения SDF-значений в разреженных воксельных блоках. Это обеспечивает кэш-дружественные (cache-friendly) запросы и ускорение рендеринга до 100 раз относительно MLP-подходов \ref{sssec:mlp}.
    \item \textbf{SDF Octrees:} Использование октодеревьев в таких методах как NGLOD \cite{takikawa2021neural}, SBS \cite{HanssonSoderlund2022SDF}, SCom Tree \cite{garifullin2025compact} позволяет реализовать адаптивный уровень детализации (Level of Detail, LoD). Например, система SCom Tree \cite{garifullin2025compact} использует кластеризацию похожих блоков данных, достигая высокой геометрической точности при экстремально низком объеме памяти и может рендериться в реальном времени используя трассировку лучей \ref{sssec:raytracing}.
\end{itemize}

\subsection{Методы рендеринга неявных поверхностей}
\label{sssec:raytracing}

В данном разделе рассматриваются современные подходы к рендерингу неявных поверхностей, с акцентом на методы, основанные на SDF.

\subsubsection{Трассировка лучей}
Трассировка лучей (Ray Tracing) \cite{hart1989ray} является классическим подходом к рендерингу неявных поверхностей. Основная идея заключается в выпуске лучей из камеры через каждый пиксель изображения и последующем поиске пересечения каждого луча с поверхностью, заданной SDF. Для эффективного поиска пересечений используются иерархические структуры данных, такие как октодеревья или BVH (Bounding Volume Hierarchy).

Алгоритм Sphere Tracing \cite{hart1989ray} является одним из наиболее распространенных методов для рендеринга SDF. Он использует свойство SDF --- гарантию того, что минимальное расстояние до поверхности в любой точке не превышает значения SDF в этой точке. Это позволяет ``прыгать'' вдоль луча на расстояние, равное значению SDF, гарантируя, что следующая точка не пересечет поверхность. Процесс повторяется до схождения или достижения максимального числа итераций.

Преимущества трассировки лучей:
\begin{itemize}
    \item Высокое качество изображения с корректными отражениями, преломлениями и глобальным освещением.
    \item Естественная поддержка сложных топологий и органических форм, характерных для SDF.
    \item Возможность использования преимуществ современных RTX-видеокарт для аппаратного ускорения.
\end{itemize}

Недостатки:
\begin{itemize}
    \item Высокая вычислительная стоимость, особенно для сложных сцен с множеством запросов SDF.
    \item Сложность реализации продвинутых оптимизаций, таких как кэширование и адаптивное сэмплирование.
    \item Проблемы с производительностью при работе с высокодетализированными сценами.
\end{itemize}

\subsubsection{Объёмный рендеринг}
\label{sssec:vrender}
Объёмный рендеринг (Volume Rendering) \cite{yariv2021volume} используется в методах, основанных на нейронных сетях (MLP), для интегрирования цвета и плотности вдоль луча. Этот подход позволяет восстанавливать не только поверхность, но и объёмные свойства сцены, что особенно полезно для полупрозрачных объектов и эффектов атмосферного рассеяния.

В контексте SDF-рендеринга объёмный рендеринг может использоваться для визуализации переходных зон вокруг поверхности, создавая эффект ``мягкой'' поверхности или гало. Однако для чисто поверхностного рендеринга более эффективными являются методы на основе прямой растеризации.

\subsection{Методы, основанные на растеризации}
В последние годы наблюдается возрождение интереса к методам растеризации для рендеринга неявных поверхностей, обусловленное развитием современных GPU-архитектур и появлением новых возможностей программируемого конвейера.

В работе Tatarchuk et al. \cite{tatarchuk2007real} впервые была продемонстрирована возможность эффективного рендеринга неявных поверхностей с использованием программируемых шейдеров. Авторы предложили подход, основанный на итеративном поиске пересечения луча с поверхностью в вершинном шейдере с последующей генерацией микрополигонов для покрытия найденной точки.

Современные подходы, такие как текущая диссертация, развивают эту идею, используя Mesh Shaders для более гибкого управления геометрией. Это позволяет генерировать треугольники непосредственно на GPU, адаптивно подстраивая их размер под видимую детализацию поверхности.

\subsection{Сравнение методов}

В таблице~\ref{tab:methods_comparison} приведено сравнение основных подходов к рендерингу SDF по ключевым характеристикам.

\input{generated/tab_methods_comparison}

\textbf{Наблюдение:} растеризация с адаптивной триангуляцией обеспечивает наилучший баланс между производительностью и качеством изображения для интерактивных приложений, особенно при использовании современных GPU с поддержкой Mesh Shaders.